\documentclass[10pt]{article}
\usepackage[T1]{fontenc}
\usepackage{lmodern}
\usepackage{amsmath,amssymb}
\usepackage[margin=0.72in]{geometry}
\usepackage[colorlinks=true,citecolor=blue,urlcolor=blue]{hyperref}

\newcommand{\Lz}{L^0}
\newcommand{\F}{\mathcal F}
\newcommand{\U}{\mathcal U}
\newcommand{\one}{\mathbf 1}

\begin{document}
\begin{center}
{\Large\bfseries Literature Status: Locally $L^0$-Convex Topologies\par}
\vspace{0.35em}
{Run \texttt{fa\_banach\_001}, agent \texttt{agent\_lane\_11}\par}
{\small 17 August 2026}
\end{center}
\vspace{0.35em}

\paragraph{Status.}
\textbf{literature\_already\_answered}.  Zapata \cite{Zapata} gives one
counterexample that answers negatively the cluster of questions following
Propositions 2.22--2.24 of Guo--Zhao--Zeng \cite{GZZ}.  We spell out two
immediate consequences of the example---its gauge is identically zero, and
no equivalent local base can consist of countable-concatenation-stable
sets---because together they match the source questions exactly.

\paragraph{The questions and the later criterion.}
For an $L^0$-balanced, absorbent, and convex set $U$ in an $L^0$-module,
write
\[
 p_U(X)=\mathop{\mathrm{ess\,inf}}\{\lambda\in L^0_{++}:X\in\lambda U\}.
\]
The source asks whether
\[
 \{X:p_U(X)<1\}\subset U                                      \tag{1}
\]
can fail without the countable concatenation property, whether every
locally $L^0$-convex topology is induced by $L^0$-seminorms, and whether
the latter conclusion can be rescued when the ambient module itself has
the countable concatenation property.

Zapata's characterization says that a topology on an $L^0$-module is
induced by a family of $L^0$-seminorms if and only if it has a zero-neighborhood
base of $L^0$-balanced, absorbent, convex sets that are closed under
countable concatenations.  His following example shows that the extra
condition is essential.

\paragraph{The counterexample.}
Let $(A_n)_{n\geq1}$ be an infinite measurable partition of $\Omega$ with
$P(A_n)>0$, let $E=L^0(\Omega,\F,P)$, and for
$\varepsilon\in L^0_{++}$ set
\[
 U_\varepsilon=
 \left\{Y\in L^0:\ \exists I\subset\mathbb N\ \text{finite such that}
 |Y\one_{A_i}|\leq\varepsilon\ \text{for every }i\notin I\right\}.       \tag{2}
\]
Zapata verifies that the family $(U_\varepsilon)$ is a zero-neighborhood
base for a locally $L^0$-convex module topology.  The elementary checks are
transparent: finite exceptional sets are preserved under balanced scaling
and finite unions handle convex combinations; absorbency follows since
\[
 \lambda=1+|Y|/\varepsilon\in L^0_{++}
 \quad\Longrightarrow\quad |Y/\lambda|\leq\varepsilon.
\]
The ambient module $E=L^0$ has the countable concatenation property.

On the other hand, every $U_\varepsilon$ is proper because
$\varepsilon+1\notin U_\varepsilon$.  Yet its countable-concatenation hull
is all of $L^0$:
\[
 H_{cc}(U_\varepsilon)=L^0.                                    \tag{3}
\]
Indeed, for any $Y\in L^0$, each one-piece truncation
$Y\one_{A_n}$ belongs to $U_\varepsilon$, and their concatenation along
$(A_n)$ is $Y$.

\paragraph{No countable-concatenation local base.}
Suppose an equivalent zero-neighborhood base contained a set $V$ closed
under countable concatenations with $V\subset U_\varepsilon$.  Because
$V$ is a neighborhood in the topology generated by (2), some
$U_\delta\subset V$.  Then (3) and concatenation stability give
\[
 L^0=H_{cc}(U_\delta)\subset V\subset U_\varepsilon,
\]
contradicting the properness of $U_\varepsilon$.  Thus even though $E$
has the countable concatenation property, its topology need not have the
kind of local base proposed in \cite{GZZ}.  By Zapata's criterion, this
topology is not induced by any family of $L^0$-seminorms.

\newpage
\paragraph{The gauge is identically zero.}
The same example gives a sharp failure of (1).  Fix $X\in L^0$, an index
$m$, and a real $t>0$.  Define
\[
 \lambda_{m,t}
 =t\one_{A_m}+\left(1+\frac{|X|}{\varepsilon}\right)
       \one_{\Omega\setminus A_m}\in L^0_{++}.
\]
Outside the single exceptional piece $A_m$,
$|X/\lambda_{m,t}|\leq\varepsilon$, so
$X/\lambda_{m,t}\in U_\varepsilon$ and hence
$p_{U_\varepsilon}(X)\leq\lambda_{m,t}$.  Restricting to $A_m$ and then
letting $t\downarrow0$ gives
$p_{U_\varepsilon}(X)=0$ on $A_m$.  Since the $A_m$ partition $\Omega$,
\[
 p_{U_\varepsilon}(X)=0\quad\text{for every }X\in L^0.          \tag{4}
\]
Consequently
\[
 \{X:p_{U_\varepsilon}(X)<1\}=L^0\not\subset U_\varepsilon,
\]
which is the requested explicit negative answer to the random-gauge
question.

\paragraph{Independent confirmation and scope.}
Guo--Zhao--Zeng later explicitly record that Wu--Guo and Zapata
independently supplied counterexamples to seminorm generation
\cite{GZZlater}.  The result above resolves the source's topology/gauge
question cluster.  It does \emph{not} claim to resolve the distinct open
problem later in \cite{GZZ}: whether an $L^0$-pre-barreled module with the
countable concatenation property must be $L^0$-barreled.  A bounded search
found no later resolution of that implication, so it is excluded from the
present result.

\begin{thebibliography}{9}
\small
\raggedright

\bibitem{GZZ}
T. Guo, S. Zhao, and X. Zeng,
\emph{On random convex analysis---the analytic foundation of the module
approach to conditional risk measures}, arXiv:1210.1848v3.

\bibitem{Zapata}
J. M. Zapata,
\emph{On characterization of locally $L^0$-convex topologies induced by a
family of $L^0$-seminorms}, arXiv:1404.0357v4.

\bibitem{GZZlater}
T. Guo, S. Zhao, and X. Zeng,
\emph{Random convex analysis (I): separation and Fenchel--Moreau duality in
random locally convex modules}, arXiv:1503.08695v2.

\end{thebibliography}

\end{document}
